
%(BEGIN_QUESTION)
% Copyright 2015, Tony R. Kuphaldt, released under the Creative Commons Attribution License (v 1.0)
% This means you may do almost anything with this work of mine, so long as you give me proper credit

Les opplæringsdokumentet "LabVIEW exercise \#1" som forberedelse til gjennomføring i datalab sammen med en USB-basert DAQ-enhet. Her lærer du å konfigurere et enkelt "virtuelt instrument" i National Instruments LabVIEW for å vise en analog signalverdi på skjermen.

\vskip 10pt

Merk: Opplæringen du trenger ligger i "Instrumentation Reference" (et sett digitale dokumenter) som instruktøren har gjort tilgjengelig.

\underbar{file i00705no}
%(END_QUESTION)





%(BEGIN_ANSWER)


%(END_ANSWER)





%(BEGIN_NOTES)

I denne øvelsen setter vi opp LabVIEW til å hente et analogt spenningssignal og vise det i Front Panel. Når oppgaven er ferdig, har vi et enkelt voltmeterdisplay på PC-en.

\vfil \eject

\noindent
{\bf Forberedelsesquiz:}

LabVIEW bruker to vinduer i utviklingsmiljøet: {\it Block Diagram} og {\it Front Panel}. Forskjellen er:

\begin{itemize}
\item{} Block Diagram viser "live" visning av alle signaler, mens Front Panel er statisk uten interaktivitet
\vskip 10pt
\item{} Block Diagram viser brukerknapper og display, mens Front Panel viser hvordan matematiske beregninger gjøres
\vskip 10pt
\item{} Block Diagram dokumenterer all elektrisk kabling på DAQ-enheten, mens Front Panel dokumenterer alle programvarefunksjoner
\vskip 10pt
\item{} Block Diagram viser signalbehandlingsfunksjonene, mens Front Panel inneholder brukerknapper og display
\end{itemize}

\vfil \eject

\noindent
{\bf Forberedelsesquiz:}

Beskriv én av aktivitetene du skal gjøre med LabVIEW og DAQ i dag. Vær så detaljert som mulig.

%INDEX% Leseoppgave: LabVIEW exercise #1

%(END_NOTES)
